% answers_short.tex —— 2022 年·新高考Ⅰ卷 · 简短答案速查
% 由 tools/make_short_answers.py 自 answers.tex 的 \gksolhead/\gkanswer 提取，
% 并与 2022/2022_数学真题_答案要点.md 逐题比对后生成，勿手改（改源 MD 后重跑）。

\begin{gknotice}
\textbf{使用说明}：本节只给最终答案便于快速核对，完整推导见解析版。答案为编写者独立推导，\textbf{非官方标准答案}。
\end{gknotice}

\gkpart{单选题}
\begin{center}\begin{tabular}{ccccccccc}
\toprule
题号 & 1 & 2 & 3 & 4 & 5 & 6 & 7 & 8 \\
\midrule
答案 & D & D & B & C & D & A & C & C \\
\bottomrule
\end{tabular}\end{center}

\gkpart{多选题}
\begin{center}\begin{tabular}{ccccc}
\toprule
题号 & 9 & 10 & 11 & 12 \\
\midrule
答案 & ABD & AC & BCD & BC \\
\bottomrule
\end{tabular}\end{center}

\gkpart{填空题}
\begin{itemize}[leftmargin=2.6em]
  \item[\textbf{13.}] $-28$
  \item[\textbf{14.}] $x=-1$（或 $3x+4y-5=0$、$7x-24y-25=0$，任写一条）
  \item[\textbf{15.}] $(-\infty,-4)\cup(0,+\infty)$
  \item[\textbf{16.}] $13$
\end{itemize}

\gkpart{解答题}
\begin{itemize}[leftmargin=2.6em]
  \item[\textbf{17.}] (1) $a_n=\dfrac{n(n+1)}2$；(2) 见解析版
  \item[\textbf{18.}] (1) $B=\dfrac\pi6$；(2) $4\sqrt2-5$
  \item[\textbf{19.}] (1) $\sqrt2$；(2) $\dfrac{\sqrt3}2$
  \item[\textbf{20.}] (1) 有 $99\%$ 把握；(2) ② $P(A\mid B)=0.4$，$P(A\mid\bar B)=0.1$，$R=6$
  \item[\textbf{21.}] (1) $k_l=-1$；(2) $\dfrac{16\sqrt2}9$
  \item[\textbf{22.}] (1) $a=1$；(2) 见解析版
\end{itemize}
